% =====================================================================
%  limitations.tex  --  PHASE VII, Task 14D
%
%  A dedicated limitations section.  It exists because the phased
%  evaluation of this framework overturned several of its own earlier
%  claims, and the resulting boundaries are specific enough to be worth
%  stating in one place rather than distributing across the results.
%
%  Intended placement: Section 11 (\label{sec:limitations}), after the
%  ablation and before the conclusion.  Also carries
%  \label{sec:discussion}.
%
%  Drop-in usage: % =====================================================================
%  limitations.tex  --  PHASE VII, Task 14D
%
%  A dedicated limitations section.  It exists because the phased
%  evaluation of this framework overturned several of its own earlier
%  claims, and the resulting boundaries are specific enough to be worth
%  stating in one place rather than distributing across the results.
%
%  Intended placement: Section 11 (\label{sec:limitations}), after the
%  ablation and before the conclusion.  Also carries
%  \label{sec:discussion}.
%
%  Drop-in usage: % =====================================================================
%  limitations.tex  --  PHASE VII, Task 14D
%
%  A dedicated limitations section.  It exists because the phased
%  evaluation of this framework overturned several of its own earlier
%  claims, and the resulting boundaries are specific enough to be worth
%  stating in one place rather than distributing across the results.
%
%  Intended placement: Section 11 (\label{sec:limitations}), after the
%  ablation and before the conclusion.  Also carries
%  \label{sec:discussion}.
%
%  Drop-in usage: % =====================================================================
%  limitations.tex  --  PHASE VII, Task 14D
%
%  A dedicated limitations section.  It exists because the phased
%  evaluation of this framework overturned several of its own earlier
%  claims, and the resulting boundaries are specific enough to be worth
%  stating in one place rather than distributing across the results.
%
%  Intended placement: Section 11 (\label{sec:limitations}), after the
%  ablation and before the conclusion.  Also carries
%  \label{sec:discussion}.
%
%  Drop-in usage: \input{limitations}
%  Bibliography: paper/references.bib
%
%  Every item below is traceable to a measured result; none is a
%  generic caveat.  Where a limitation is quantified, the number is the
%  one reported in the section named.
% =====================================================================

\section{Limitations}
\label{sec:limitations}
\label{sec:discussion}

The evaluation in this paper was run in phases, each testing the
preceding one, and several of the framework's original claims did not
survive that process. The boundaries this leaves are specific, and are
collected here rather than distributed through the results.

\subsection{The economic claim rests on eighteen factory days}
\label{sec:limitations:sample}

The decision layer is evaluated on $394$ held-out idle episodes spread
over \textbf{18 factory days}. The day-block bootstrap therefore has
eighteen units to resample, and no refinement of the estimator narrows
the resulting intervals: the proposed policy's $+2.9$\,USD/yr sits
inside $[-55.9,+46.2]$, and every policy's interval against the plant's
own behaviour spans zero except the oracle's
(Section~\ref{sec:significance:decision}).

\textbf{We therefore do not claim that the proposed policy saves money
against the status quo.} What the data support is narrower and is stated
as such throughout: adaptive thresholds beat the fixed threshold they
replace ($+68.9$\,USD/yr, $p=0.001$, paired on the same resample), and
the head-room is real but largely uncaptured. Whether forecasting beats
not forecasting on this machine is an open question that eighteen days
cannot answer, and the ten-seed spread of $\pm0.6$\,USD/yr against a
sampling spread of $\pm50$\,USD/yr confirms that the binding constraint
is the record, not the model.

A second consequence follows for anyone reproducing this work: the
scenario axis is not the uncertainty. The three scenarios differ in the
valuation of lost production, which is the one coefficient that could
not be measured and is therefore reported as an \emph{output} of the
break-even identity rather than assumed. A plant that supplies its own
figure collapses the three columns to one; it does not narrow the
interval.

\subsection{Two machines cannot be annualised, and the largest
per-day figure is one of them}
\label{sec:limitations:milling}

The milling machines' records cover $12.2$ of $43$ days, in $1\,082$ and
$580$ acquisition fragments. Twelve covered days support two one-week
bootstrap blocks, which is not enough to form a usable interval; the
intervals printed for those rows in Table~\ref{tab:multimachine} exist
only to make that visible.

This matters because their $5.5$\,h/day is the \emph{largest} per-day
standby figure among the four machines on which the state definition is
physically valid, and it is exactly the number a reader would want to
annualise. \textbf{It must not be annualised.} The same fragmentation
also breaks the decision layer on those machines --- forecast MAE
$18\,200$\,s against $9\,034$\,s on pelletizer~I, and a proposed policy
that loses over $1\,000$\,USD/yr
(Section~\ref{sec:multimachine:decision}) --- so the two limitations have
one cause.

\subsection{The window forecaster leans on a site-specific quantity}
\label{sec:limitations:voltage}

Supply voltage is the largest single channel contributor to the window
forecaster: $20.0\%$ of its Shapley attribution across nine design
columns, $11.8\%$ on the final observed value alone
(Section~\ref{sec:explain:window}). Voltage is not a state variable of
the machine. It sags under load and recovers when the drive unloads, so
it is a real but indirect indicator --- and one determined by this site's
point of common coupling, transformer impedance and the other loads
sharing the bus.

\textbf{A deployment on another site should not expect that contribution
to transfer}, and the zero-shot cross-site failure of
Section~\ref{sec:crossmachine:dataset} (Spearman $-0.245$, worse than
chance ordering) is consistent with it. Split gain ranks voltage tenth
of twenty-five channels and would have given no warning; the limitation
was found by attribution and would not otherwise have been stated at
all.

\subsection{The attributions are single-machine, and one step removed}
\label{sec:limitations:attribution}

Three separate qualifications apply to Section~\ref{sec:explain}.

\emph{One machine.} The transition structure is measured on all eight
machines, but the forecaster and state-layer attributions are computed
on pelletizer~I only. Whether power factor leads the STANDBY split on
pelletizer~II and the milling machines is a cheap experiment that has
not been run, and it is the obvious strengthening of the physics claim.

\emph{A surrogate is not the model.} The state layer is a GMM--HMM;
what is attributed is a gradient-boosted tree fitted to its output. The
rankings are meaningful and agree with two other measures, but the
mechanism is one step removed, and no attribution method for a
Viterbi-decoded mixture model is applied here.

\emph{Feature dependence is not modelled.} Tree SHAP with the
path-dependent perturbation splits credit between correlated features in
a way that depends on tree structure, and \texttt{active\_power},
\texttt{apparent\_power} and \texttt{current} are strongly correlated on
this record. The ordering \emph{among those three} should be treated as
less firm than the gap between them and the calendar or
rolling-statistic groups.

\subsection{What the ablation does not cover}
\label{sec:limitations:ablation}

Table~\ref{tab:ablation} is computed on one machine and, in its decision
block, on one scenario (S2). The other two scenarios are stored and give
the same ordering, but they are not in the table.

More importantly, \textbf{the strongest accept/reject result in this
paper is not a row of the ablation}. Conditioning the labelling on a
passing physics score collapses the STANDBY-hour spread from $\pm33.7\%$
to $\pm3.0\%$ (Section~\ref{sec:significance:seeds}) --- a larger effect
than any component the table does contain. It is absent because it is
measured over labelling seeds while every other row is measured over
held-out blocks, windows or days, and putting incommensurable units in
one table is precisely what Section~\ref{sec:ablation:structure} argues
against. A seed-resampled ablation of the battery, check by check, is
the experiment that would close this gap.

\subsection{Hyperparameters that are stated rather than learned}
\label{sec:limitations:hyper}

Four constants are set by argument and not by search, and each is a
place where a different plant would need to re-derive rather than
re-use.

The \emph{minimum dwell} is fixed at $10$\,s, equal to the flicker
threshold, so that it removes exactly what the metric defines as
implausible and nothing more. That avoids tuning it to its own metric
but does not make it learned; a hidden semi-Markov
model~\cite{yu2010hsmm} would estimate the duration distribution
instead. The \emph{minimum off time} is $600$\,s, with the saving flat
between $300$ and $1\,200$\,s. The \emph{chance-constraint confidence}
is a point value of $0.90$, stated before the sweep that later showed it
near-optimal; a full stochastic programme would price infeasibility
rather than constrain its probability. And $k=4$ is \emph{imposed}: the
auxiliary machines would be better described with two states and the
BIC scan says the driven machines want more than four
(Section~\ref{sec:multimachine:general}), so a per-machine $k$ selected
under the physics veto is a needed extension.

One further modelling gap belongs here. \textbf{Restart wear is priced
only through a per-day cardinality cap}, not as a cost per cycle with a
bearing-life model --- and the sensitivity sweep found that tightening
that cap to one restart per day \emph{more than doubles} the saving
($89$ against $39$\,USD/yr). Since the cap is applied chronologically it
cannot be selecting the best shutdowns, so a blunt restriction improving
the outcome means the policy's marginal shutdowns are loss-making: the
greedy online rule is not optimal under a binding cap, and a positive
expected-saving margin rather than a zero one is the obvious refinement.
It is reported rather than tuned away.

\subsection{Scope of the evidence}
\label{sec:limitations:scope}

\emph{One facility.} All eight machines share a plant, a shift calendar
and an acquisition system --- and the plant calendar is precisely what
the duration forecaster turns out to be using. Cross-facility transfer
of that component is untested on a comparable record, and the
single-channel evidence available suggests it would fail.

\emph{One comparable dataset, and it is single-channel.} The second
dataset can exercise the pipeline but not the physics battery: C3 and
C4, the two checks that do not depend on power magnitude, are undefined
without current and power-factor channels, and they are the two that
carry the discriminating power
(Section~\ref{sec:crossmachine:dataset}).
On a single-channel record the framework still produces states and still
consumes them, and the negative control shows it will report recoverable
standby on a device that cannot be shut down. A second \emph{six-channel}
industrial record would test the part of the framework that matters
most; none was available.

\emph{No contactor ground truth.} Nothing in this work observes the
machine's contactor directly, so the STANDBY/OFF boundary is inferred
throughout and its residual uncertainty is irreducible from these
measurements alone. Every state-time total is accordingly quoted to two
significant figures and conditional on the physics battery passing.

\subsection{Where the framework should not be deployed}
\label{sec:limitations:deploy}

Stated as criteria rather than caveats, the evidence in this paper rules
out four situations. \textbf{Non-motor loads}: the state definition is
specific to induction-motor drives, and on the fans and contactors there
is no energised-idle state to find. \textbf{Sub-watt standby power}: the
break-even point is then numerically undefined and every currency figure
is an artefact. \textbf{Heavily fragmented records}: the decision layer
inherits its forecaster's quality, and the milling machines show what
that costs. \textbf{Single-channel metering}: the validation battery
loses the two checks that carry its discriminating power. In the first
two cases the framework detects the condition itself and declines to
report a decision, which is the intended behaviour; in the last two it
does not, and an external check is required.

%  Bibliography: paper/references.bib
%
%  Every item below is traceable to a measured result; none is a
%  generic caveat.  Where a limitation is quantified, the number is the
%  one reported in the section named.
% =====================================================================

\section{Limitations}
\label{sec:limitations}
\label{sec:discussion}

The evaluation in this paper was run in phases, each testing the
preceding one, and several of the framework's original claims did not
survive that process. The boundaries this leaves are specific, and are
collected here rather than distributed through the results.

\subsection{The economic claim rests on eighteen factory days}
\label{sec:limitations:sample}

The decision layer is evaluated on $394$ held-out idle episodes spread
over \textbf{18 factory days}. The day-block bootstrap therefore has
eighteen units to resample, and no refinement of the estimator narrows
the resulting intervals: the proposed policy's $+2.9$\,USD/yr sits
inside $[-55.9,+46.2]$, and every policy's interval against the plant's
own behaviour spans zero except the oracle's
(Section~\ref{sec:significance:decision}).

\textbf{We therefore do not claim that the proposed policy saves money
against the status quo.} What the data support is narrower and is stated
as such throughout: adaptive thresholds beat the fixed threshold they
replace ($+68.9$\,USD/yr, $p=0.001$, paired on the same resample), and
the head-room is real but largely uncaptured. Whether forecasting beats
not forecasting on this machine is an open question that eighteen days
cannot answer, and the ten-seed spread of $\pm0.6$\,USD/yr against a
sampling spread of $\pm50$\,USD/yr confirms that the binding constraint
is the record, not the model.

A second consequence follows for anyone reproducing this work: the
scenario axis is not the uncertainty. The three scenarios differ in the
valuation of lost production, which is the one coefficient that could
not be measured and is therefore reported as an \emph{output} of the
break-even identity rather than assumed. A plant that supplies its own
figure collapses the three columns to one; it does not narrow the
interval.

\subsection{Two machines cannot be annualised, and the largest
per-day figure is one of them}
\label{sec:limitations:milling}

The milling machines' records cover $12.2$ of $43$ days, in $1\,082$ and
$580$ acquisition fragments. Twelve covered days support two one-week
bootstrap blocks, which is not enough to form a usable interval; the
intervals printed for those rows in Table~\ref{tab:multimachine} exist
only to make that visible.

This matters because their $5.5$\,h/day is the \emph{largest} per-day
standby figure among the four machines on which the state definition is
physically valid, and it is exactly the number a reader would want to
annualise. \textbf{It must not be annualised.} The same fragmentation
also breaks the decision layer on those machines --- forecast MAE
$18\,200$\,s against $9\,034$\,s on pelletizer~I, and a proposed policy
that loses over $1\,000$\,USD/yr
(Section~\ref{sec:multimachine:decision}) --- so the two limitations have
one cause.

\subsection{The window forecaster leans on a site-specific quantity}
\label{sec:limitations:voltage}

Supply voltage is the largest single channel contributor to the window
forecaster: $20.0\%$ of its Shapley attribution across nine design
columns, $11.8\%$ on the final observed value alone
(Section~\ref{sec:explain:window}). Voltage is not a state variable of
the machine. It sags under load and recovers when the drive unloads, so
it is a real but indirect indicator --- and one determined by this site's
point of common coupling, transformer impedance and the other loads
sharing the bus.

\textbf{A deployment on another site should not expect that contribution
to transfer}, and the zero-shot cross-site failure of
Section~\ref{sec:crossmachine:dataset} (Spearman $-0.245$, worse than
chance ordering) is consistent with it. Split gain ranks voltage tenth
of twenty-five channels and would have given no warning; the limitation
was found by attribution and would not otherwise have been stated at
all.

\subsection{The attributions are single-machine, and one step removed}
\label{sec:limitations:attribution}

Three separate qualifications apply to Section~\ref{sec:explain}.

\emph{One machine.} The transition structure is measured on all eight
machines, but the forecaster and state-layer attributions are computed
on pelletizer~I only. Whether power factor leads the STANDBY split on
pelletizer~II and the milling machines is a cheap experiment that has
not been run, and it is the obvious strengthening of the physics claim.

\emph{A surrogate is not the model.} The state layer is a GMM--HMM;
what is attributed is a gradient-boosted tree fitted to its output. The
rankings are meaningful and agree with two other measures, but the
mechanism is one step removed, and no attribution method for a
Viterbi-decoded mixture model is applied here.

\emph{Feature dependence is not modelled.} Tree SHAP with the
path-dependent perturbation splits credit between correlated features in
a way that depends on tree structure, and \texttt{active\_power},
\texttt{apparent\_power} and \texttt{current} are strongly correlated on
this record. The ordering \emph{among those three} should be treated as
less firm than the gap between them and the calendar or
rolling-statistic groups.

\subsection{What the ablation does not cover}
\label{sec:limitations:ablation}

Table~\ref{tab:ablation} is computed on one machine and, in its decision
block, on one scenario (S2). The other two scenarios are stored and give
the same ordering, but they are not in the table.

More importantly, \textbf{the strongest accept/reject result in this
paper is not a row of the ablation}. Conditioning the labelling on a
passing physics score collapses the STANDBY-hour spread from $\pm33.7\%$
to $\pm3.0\%$ (Section~\ref{sec:significance:seeds}) --- a larger effect
than any component the table does contain. It is absent because it is
measured over labelling seeds while every other row is measured over
held-out blocks, windows or days, and putting incommensurable units in
one table is precisely what Section~\ref{sec:ablation:structure} argues
against. A seed-resampled ablation of the battery, check by check, is
the experiment that would close this gap.

\subsection{Hyperparameters that are stated rather than learned}
\label{sec:limitations:hyper}

Four constants are set by argument and not by search, and each is a
place where a different plant would need to re-derive rather than
re-use.

The \emph{minimum dwell} is fixed at $10$\,s, equal to the flicker
threshold, so that it removes exactly what the metric defines as
implausible and nothing more. That avoids tuning it to its own metric
but does not make it learned; a hidden semi-Markov
model~\cite{yu2010hsmm} would estimate the duration distribution
instead. The \emph{minimum off time} is $600$\,s, with the saving flat
between $300$ and $1\,200$\,s. The \emph{chance-constraint confidence}
is a point value of $0.90$, stated before the sweep that later showed it
near-optimal; a full stochastic programme would price infeasibility
rather than constrain its probability. And $k=4$ is \emph{imposed}: the
auxiliary machines would be better described with two states and the
BIC scan says the driven machines want more than four
(Section~\ref{sec:multimachine:general}), so a per-machine $k$ selected
under the physics veto is a needed extension.

One further modelling gap belongs here. \textbf{Restart wear is priced
only through a per-day cardinality cap}, not as a cost per cycle with a
bearing-life model --- and the sensitivity sweep found that tightening
that cap to one restart per day \emph{more than doubles} the saving
($89$ against $39$\,USD/yr). Since the cap is applied chronologically it
cannot be selecting the best shutdowns, so a blunt restriction improving
the outcome means the policy's marginal shutdowns are loss-making: the
greedy online rule is not optimal under a binding cap, and a positive
expected-saving margin rather than a zero one is the obvious refinement.
It is reported rather than tuned away.

\subsection{Scope of the evidence}
\label{sec:limitations:scope}

\emph{One facility.} All eight machines share a plant, a shift calendar
and an acquisition system --- and the plant calendar is precisely what
the duration forecaster turns out to be using. Cross-facility transfer
of that component is untested on a comparable record, and the
single-channel evidence available suggests it would fail.

\emph{One comparable dataset, and it is single-channel.} The second
dataset can exercise the pipeline but not the physics battery: C3 and
C4, the two checks that do not depend on power magnitude, are undefined
without current and power-factor channels, and they are the two that
carry the discriminating power
(Section~\ref{sec:crossmachine:dataset}).
On a single-channel record the framework still produces states and still
consumes them, and the negative control shows it will report recoverable
standby on a device that cannot be shut down. A second \emph{six-channel}
industrial record would test the part of the framework that matters
most; none was available.

\emph{No contactor ground truth.} Nothing in this work observes the
machine's contactor directly, so the STANDBY/OFF boundary is inferred
throughout and its residual uncertainty is irreducible from these
measurements alone. Every state-time total is accordingly quoted to two
significant figures and conditional on the physics battery passing.

\subsection{Where the framework should not be deployed}
\label{sec:limitations:deploy}

Stated as criteria rather than caveats, the evidence in this paper rules
out four situations. \textbf{Non-motor loads}: the state definition is
specific to induction-motor drives, and on the fans and contactors there
is no energised-idle state to find. \textbf{Sub-watt standby power}: the
break-even point is then numerically undefined and every currency figure
is an artefact. \textbf{Heavily fragmented records}: the decision layer
inherits its forecaster's quality, and the milling machines show what
that costs. \textbf{Single-channel metering}: the validation battery
loses the two checks that carry its discriminating power. In the first
two cases the framework detects the condition itself and declines to
report a decision, which is the intended behaviour; in the last two it
does not, and an external check is required.

%  Bibliography: paper/references.bib
%
%  Every item below is traceable to a measured result; none is a
%  generic caveat.  Where a limitation is quantified, the number is the
%  one reported in the section named.
% =====================================================================

\section{Limitations}
\label{sec:limitations}
\label{sec:discussion}

The evaluation in this paper was run in phases, each testing the
preceding one, and several of the framework's original claims did not
survive that process. The boundaries this leaves are specific, and are
collected here rather than distributed through the results.

\subsection{The economic claim rests on eighteen factory days}
\label{sec:limitations:sample}

The decision layer is evaluated on $394$ held-out idle episodes spread
over \textbf{18 factory days}. The day-block bootstrap therefore has
eighteen units to resample, and no refinement of the estimator narrows
the resulting intervals: the proposed policy's $+2.9$\,USD/yr sits
inside $[-55.9,+46.2]$, and every policy's interval against the plant's
own behaviour spans zero except the oracle's
(Section~\ref{sec:significance:decision}).

\textbf{We therefore do not claim that the proposed policy saves money
against the status quo.} What the data support is narrower and is stated
as such throughout: adaptive thresholds beat the fixed threshold they
replace ($+68.9$\,USD/yr, $p=0.001$, paired on the same resample), and
the head-room is real but largely uncaptured. Whether forecasting beats
not forecasting on this machine is an open question that eighteen days
cannot answer, and the ten-seed spread of $\pm0.6$\,USD/yr against a
sampling spread of $\pm50$\,USD/yr confirms that the binding constraint
is the record, not the model.

A second consequence follows for anyone reproducing this work: the
scenario axis is not the uncertainty. The three scenarios differ in the
valuation of lost production, which is the one coefficient that could
not be measured and is therefore reported as an \emph{output} of the
break-even identity rather than assumed. A plant that supplies its own
figure collapses the three columns to one; it does not narrow the
interval.

\subsection{Two machines cannot be annualised, and the largest
per-day figure is one of them}
\label{sec:limitations:milling}

The milling machines' records cover $12.2$ of $43$ days, in $1\,082$ and
$580$ acquisition fragments. Twelve covered days support two one-week
bootstrap blocks, which is not enough to form a usable interval; the
intervals printed for those rows in Table~\ref{tab:multimachine} exist
only to make that visible.

This matters because their $5.5$\,h/day is the \emph{largest} per-day
standby figure among the four machines on which the state definition is
physically valid, and it is exactly the number a reader would want to
annualise. \textbf{It must not be annualised.} The same fragmentation
also breaks the decision layer on those machines --- forecast MAE
$18\,200$\,s against $9\,034$\,s on pelletizer~I, and a proposed policy
that loses over $1\,000$\,USD/yr
(Section~\ref{sec:multimachine:decision}) --- so the two limitations have
one cause.

\subsection{The window forecaster leans on a site-specific quantity}
\label{sec:limitations:voltage}

Supply voltage is the largest single channel contributor to the window
forecaster: $20.0\%$ of its Shapley attribution across nine design
columns, $11.8\%$ on the final observed value alone
(Section~\ref{sec:explain:window}). Voltage is not a state variable of
the machine. It sags under load and recovers when the drive unloads, so
it is a real but indirect indicator --- and one determined by this site's
point of common coupling, transformer impedance and the other loads
sharing the bus.

\textbf{A deployment on another site should not expect that contribution
to transfer}, and the zero-shot cross-site failure of
Section~\ref{sec:crossmachine:dataset} (Spearman $-0.245$, worse than
chance ordering) is consistent with it. Split gain ranks voltage tenth
of twenty-five channels and would have given no warning; the limitation
was found by attribution and would not otherwise have been stated at
all.

\subsection{The attributions are single-machine, and one step removed}
\label{sec:limitations:attribution}

Three separate qualifications apply to Section~\ref{sec:explain}.

\emph{One machine.} The transition structure is measured on all eight
machines, but the forecaster and state-layer attributions are computed
on pelletizer~I only. Whether power factor leads the STANDBY split on
pelletizer~II and the milling machines is a cheap experiment that has
not been run, and it is the obvious strengthening of the physics claim.

\emph{A surrogate is not the model.} The state layer is a GMM--HMM;
what is attributed is a gradient-boosted tree fitted to its output. The
rankings are meaningful and agree with two other measures, but the
mechanism is one step removed, and no attribution method for a
Viterbi-decoded mixture model is applied here.

\emph{Feature dependence is not modelled.} Tree SHAP with the
path-dependent perturbation splits credit between correlated features in
a way that depends on tree structure, and \texttt{active\_power},
\texttt{apparent\_power} and \texttt{current} are strongly correlated on
this record. The ordering \emph{among those three} should be treated as
less firm than the gap between them and the calendar or
rolling-statistic groups.

\subsection{What the ablation does not cover}
\label{sec:limitations:ablation}

Table~\ref{tab:ablation} is computed on one machine and, in its decision
block, on one scenario (S2). The other two scenarios are stored and give
the same ordering, but they are not in the table.

More importantly, \textbf{the strongest accept/reject result in this
paper is not a row of the ablation}. Conditioning the labelling on a
passing physics score collapses the STANDBY-hour spread from $\pm33.7\%$
to $\pm3.0\%$ (Section~\ref{sec:significance:seeds}) --- a larger effect
than any component the table does contain. It is absent because it is
measured over labelling seeds while every other row is measured over
held-out blocks, windows or days, and putting incommensurable units in
one table is precisely what Section~\ref{sec:ablation:structure} argues
against. A seed-resampled ablation of the battery, check by check, is
the experiment that would close this gap.

\subsection{Hyperparameters that are stated rather than learned}
\label{sec:limitations:hyper}

Four constants are set by argument and not by search, and each is a
place where a different plant would need to re-derive rather than
re-use.

The \emph{minimum dwell} is fixed at $10$\,s, equal to the flicker
threshold, so that it removes exactly what the metric defines as
implausible and nothing more. That avoids tuning it to its own metric
but does not make it learned; a hidden semi-Markov
model~\cite{yu2010hsmm} would estimate the duration distribution
instead. The \emph{minimum off time} is $600$\,s, with the saving flat
between $300$ and $1\,200$\,s. The \emph{chance-constraint confidence}
is a point value of $0.90$, stated before the sweep that later showed it
near-optimal; a full stochastic programme would price infeasibility
rather than constrain its probability. And $k=4$ is \emph{imposed}: the
auxiliary machines would be better described with two states and the
BIC scan says the driven machines want more than four
(Section~\ref{sec:multimachine:general}), so a per-machine $k$ selected
under the physics veto is a needed extension.

One further modelling gap belongs here. \textbf{Restart wear is priced
only through a per-day cardinality cap}, not as a cost per cycle with a
bearing-life model --- and the sensitivity sweep found that tightening
that cap to one restart per day \emph{more than doubles} the saving
($89$ against $39$\,USD/yr). Since the cap is applied chronologically it
cannot be selecting the best shutdowns, so a blunt restriction improving
the outcome means the policy's marginal shutdowns are loss-making: the
greedy online rule is not optimal under a binding cap, and a positive
expected-saving margin rather than a zero one is the obvious refinement.
It is reported rather than tuned away.

\subsection{Scope of the evidence}
\label{sec:limitations:scope}

\emph{One facility.} All eight machines share a plant, a shift calendar
and an acquisition system --- and the plant calendar is precisely what
the duration forecaster turns out to be using. Cross-facility transfer
of that component is untested on a comparable record, and the
single-channel evidence available suggests it would fail.

\emph{One comparable dataset, and it is single-channel.} The second
dataset can exercise the pipeline but not the physics battery: C3 and
C4, the two checks that do not depend on power magnitude, are undefined
without current and power-factor channels, and they are the two that
carry the discriminating power
(Section~\ref{sec:crossmachine:dataset}).
On a single-channel record the framework still produces states and still
consumes them, and the negative control shows it will report recoverable
standby on a device that cannot be shut down. A second \emph{six-channel}
industrial record would test the part of the framework that matters
most; none was available.

\emph{No contactor ground truth.} Nothing in this work observes the
machine's contactor directly, so the STANDBY/OFF boundary is inferred
throughout and its residual uncertainty is irreducible from these
measurements alone. Every state-time total is accordingly quoted to two
significant figures and conditional on the physics battery passing.

\subsection{Where the framework should not be deployed}
\label{sec:limitations:deploy}

Stated as criteria rather than caveats, the evidence in this paper rules
out four situations. \textbf{Non-motor loads}: the state definition is
specific to induction-motor drives, and on the fans and contactors there
is no energised-idle state to find. \textbf{Sub-watt standby power}: the
break-even point is then numerically undefined and every currency figure
is an artefact. \textbf{Heavily fragmented records}: the decision layer
inherits its forecaster's quality, and the milling machines show what
that costs. \textbf{Single-channel metering}: the validation battery
loses the two checks that carry its discriminating power. In the first
two cases the framework detects the condition itself and declines to
report a decision, which is the intended behaviour; in the last two it
does not, and an external check is required.

%  Bibliography: paper/references.bib
%
%  Every item below is traceable to a measured result; none is a
%  generic caveat.  Where a limitation is quantified, the number is the
%  one reported in the section named.
% =====================================================================

\section{Limitations}
\label{sec:limitations}
\label{sec:discussion}

The evaluation in this paper was run in phases, each testing the
preceding one, and several of the framework's original claims did not
survive that process. The boundaries this leaves are specific, and are
collected here rather than distributed through the results.

\subsection{The economic claim rests on eighteen factory days}
\label{sec:limitations:sample}

The decision layer is evaluated on $394$ held-out idle episodes spread
over \textbf{18 factory days}. The day-block bootstrap therefore has
eighteen units to resample, and no refinement of the estimator narrows
the resulting intervals: the proposed policy's $+2.9$\,USD/yr sits
inside $[-55.9,+46.2]$, and every policy's interval against the plant's
own behaviour spans zero except the oracle's
(Section~\ref{sec:significance:decision}).

\textbf{We therefore do not claim that the proposed policy saves money
against the status quo.} What the data support is narrower and is stated
as such throughout: adaptive thresholds beat the fixed threshold they
replace ($+68.9$\,USD/yr, $p=0.001$, paired on the same resample), and
the head-room is real but largely uncaptured. Whether forecasting beats
not forecasting on this machine is an open question that eighteen days
cannot answer, and the ten-seed spread of $\pm0.6$\,USD/yr against a
sampling spread of $\pm50$\,USD/yr confirms that the binding constraint
is the record, not the model.

A second consequence follows for anyone reproducing this work: the
scenario axis is not the uncertainty. The three scenarios differ in the
valuation of lost production, which is the one coefficient that could
not be measured and is therefore reported as an \emph{output} of the
break-even identity rather than assumed. A plant that supplies its own
figure collapses the three columns to one; it does not narrow the
interval.

\subsection{Two machines cannot be annualised, and the largest
per-day figure is one of them}
\label{sec:limitations:milling}

The milling machines' records cover $12.2$ of $43$ days, in $1\,082$ and
$580$ acquisition fragments. Twelve covered days support two one-week
bootstrap blocks, which is not enough to form a usable interval; the
intervals printed for those rows in Table~\ref{tab:multimachine} exist
only to make that visible.

This matters because their $5.5$\,h/day is the \emph{largest} per-day
standby figure among the four machines on which the state definition is
physically valid, and it is exactly the number a reader would want to
annualise. \textbf{It must not be annualised.} The same fragmentation
also breaks the decision layer on those machines --- forecast MAE
$18\,200$\,s against $9\,034$\,s on pelletizer~I, and a proposed policy
that loses over $1\,000$\,USD/yr
(Section~\ref{sec:multimachine:decision}) --- so the two limitations have
one cause.

\subsection{The window forecaster leans on a site-specific quantity}
\label{sec:limitations:voltage}

Supply voltage is the largest single channel contributor to the window
forecaster: $20.0\%$ of its Shapley attribution across nine design
columns, $11.8\%$ on the final observed value alone
(Section~\ref{sec:explain:window}). Voltage is not a state variable of
the machine. It sags under load and recovers when the drive unloads, so
it is a real but indirect indicator --- and one determined by this site's
point of common coupling, transformer impedance and the other loads
sharing the bus.

\textbf{A deployment on another site should not expect that contribution
to transfer}, and the zero-shot cross-site failure of
Section~\ref{sec:crossmachine:dataset} (Spearman $-0.245$, worse than
chance ordering) is consistent with it. Split gain ranks voltage tenth
of twenty-five channels and would have given no warning; the limitation
was found by attribution and would not otherwise have been stated at
all.

\subsection{The attributions are single-machine, and one step removed}
\label{sec:limitations:attribution}

Three separate qualifications apply to Section~\ref{sec:explain}.

\emph{One machine.} The transition structure is measured on all eight
machines, but the forecaster and state-layer attributions are computed
on pelletizer~I only. Whether power factor leads the STANDBY split on
pelletizer~II and the milling machines is a cheap experiment that has
not been run, and it is the obvious strengthening of the physics claim.

\emph{A surrogate is not the model.} The state layer is a GMM--HMM;
what is attributed is a gradient-boosted tree fitted to its output. The
rankings are meaningful and agree with two other measures, but the
mechanism is one step removed, and no attribution method for a
Viterbi-decoded mixture model is applied here.

\emph{Feature dependence is not modelled.} Tree SHAP with the
path-dependent perturbation splits credit between correlated features in
a way that depends on tree structure, and \texttt{active\_power},
\texttt{apparent\_power} and \texttt{current} are strongly correlated on
this record. The ordering \emph{among those three} should be treated as
less firm than the gap between them and the calendar or
rolling-statistic groups.

\subsection{What the ablation does not cover}
\label{sec:limitations:ablation}

Table~\ref{tab:ablation} is computed on one machine and, in its decision
block, on one scenario (S2). The other two scenarios are stored and give
the same ordering, but they are not in the table.

More importantly, \textbf{the strongest accept/reject result in this
paper is not a row of the ablation}. Conditioning the labelling on a
passing physics score collapses the STANDBY-hour spread from $\pm33.7\%$
to $\pm3.0\%$ (Section~\ref{sec:significance:seeds}) --- a larger effect
than any component the table does contain. It is absent because it is
measured over labelling seeds while every other row is measured over
held-out blocks, windows or days, and putting incommensurable units in
one table is precisely what Section~\ref{sec:ablation:structure} argues
against. A seed-resampled ablation of the battery, check by check, is
the experiment that would close this gap.

\subsection{Hyperparameters that are stated rather than learned}
\label{sec:limitations:hyper}

Four constants are set by argument and not by search, and each is a
place where a different plant would need to re-derive rather than
re-use.

The \emph{minimum dwell} is fixed at $10$\,s, equal to the flicker
threshold, so that it removes exactly what the metric defines as
implausible and nothing more. That avoids tuning it to its own metric
but does not make it learned; a hidden semi-Markov
model~\cite{yu2010hsmm} would estimate the duration distribution
instead. The \emph{minimum off time} is $600$\,s, with the saving flat
between $300$ and $1\,200$\,s. The \emph{chance-constraint confidence}
is a point value of $0.90$, stated before the sweep that later showed it
near-optimal; a full stochastic programme would price infeasibility
rather than constrain its probability. And $k=4$ is \emph{imposed}: the
auxiliary machines would be better described with two states and the
BIC scan says the driven machines want more than four
(Section~\ref{sec:multimachine:general}), so a per-machine $k$ selected
under the physics veto is a needed extension.

One further modelling gap belongs here. \textbf{Restart wear is priced
only through a per-day cardinality cap}, not as a cost per cycle with a
bearing-life model --- and the sensitivity sweep found that tightening
that cap to one restart per day \emph{more than doubles} the saving
($89$ against $39$\,USD/yr). Since the cap is applied chronologically it
cannot be selecting the best shutdowns, so a blunt restriction improving
the outcome means the policy's marginal shutdowns are loss-making: the
greedy online rule is not optimal under a binding cap, and a positive
expected-saving margin rather than a zero one is the obvious refinement.
It is reported rather than tuned away.

\subsection{Scope of the evidence}
\label{sec:limitations:scope}

\emph{One facility.} All eight machines share a plant, a shift calendar
and an acquisition system --- and the plant calendar is precisely what
the duration forecaster turns out to be using. Cross-facility transfer
of that component is untested on a comparable record, and the
single-channel evidence available suggests it would fail.

\emph{One comparable dataset, and it is single-channel.} The second
dataset can exercise the pipeline but not the physics battery: C3 and
C4, the two checks that do not depend on power magnitude, are undefined
without current and power-factor channels, and they are the two that
carry the discriminating power
(Section~\ref{sec:crossmachine:dataset}).
On a single-channel record the framework still produces states and still
consumes them, and the negative control shows it will report recoverable
standby on a device that cannot be shut down. A second \emph{six-channel}
industrial record would test the part of the framework that matters
most; none was available.

\emph{No contactor ground truth.} Nothing in this work observes the
machine's contactor directly, so the STANDBY/OFF boundary is inferred
throughout and its residual uncertainty is irreducible from these
measurements alone. Every state-time total is accordingly quoted to two
significant figures and conditional on the physics battery passing.

\subsection{Where the framework should not be deployed}
\label{sec:limitations:deploy}

Stated as criteria rather than caveats, the evidence in this paper rules
out four situations. \textbf{Non-motor loads}: the state definition is
specific to induction-motor drives, and on the fans and contactors there
is no energised-idle state to find. \textbf{Sub-watt standby power}: the
break-even point is then numerically undefined and every currency figure
is an artefact. \textbf{Heavily fragmented records}: the decision layer
inherits its forecaster's quality, and the milling machines show what
that costs. \textbf{Single-channel metering}: the validation battery
loses the two checks that carry its discriminating power. In the first
two cases the framework detects the condition itself and declines to
report a decision, which is the intended behaviour; in the last two it
does not, and an external check is required.
